\chapter{Introduction}
\label{ch:intro}

\section{Motivation}
Programming assessments increasingly happen on a computer, in a lab or
remotely. Identity is checked when a session starts, and the final code is
checked for copying or generated text when it ends. A student can pass both
checks while somebody else wrote most of the code: a friend who takes the
keyboard for twenty minutes, or a paid proxy on a remote call.
\todo{Evidence of the problem: contract cheating / proxy statistics.}

\section{Authorship continuity}
\begin{definition}[Authorship continuity]
A coding session has authorship continuity if the same person produced all
of its keystrokes and edits.
\end{definition}
This thesis asks whether a break in continuity can be detected from how
code is written (process), without watching the student.

\section{Research questions}
\begin{description}[leftmargin=2.5em]
  \item[\RQ{1}] Can a mid-session change of author be detected, and how
    quickly and with how many false alarms? \emph{(main contribution)}
  \item[\RQ{2}] Can composing original code be distinguished from
    transcribing code from another source?
  \item[\RQ{3}] Does a baseline learned in the lab transfer to a remote
    interview on the student's own laptop?
  \item[\RQ{4}] Is the detector fair across subgroups, and what is its
    precision at realistic cheating rates (1 in 100, 1 in 500)?
\end{description}

\section{Contributions}
\begin{enumerate}
  \item The Student IDE, a privacy-preserving event-logging editor.
  \item Per-student habit baselines over 30-second windows.
  \item A CUSUM-based switch detector that also estimates the switch time.
  \item A teacher-facing evidence timeline designed so a flag leads to a
    conversation, not a penalty.
  \item A staged-takeover dataset and evaluation with fairness and
    prevalence-adjusted precision.
\end{enumerate}

\section{Thesis outline}
\todo{One line per chapter once chapters stabilise.}
